\SourcePage{14}
原先曾计划由O. Gabber共同署名本文。他为了不对文中的错误或不够准确之处承担共同责任，决定不列名为作者。不过，我们采用的许多思想都出自他；他也提出了大量批评意见，我们希望这些意见已经使手稿有所改进，读者对此应当感谢他。

M. Goresky与R. MacPherson在[4]、[5]、[6]中定义并研究了适当空间$X$的\term{intersection-cohomology}{交叉上同调}。在[4]、[5]中，$X$是一个定向PL伪流形。也就是说，若$X$的维数为$n$，则它是一个PL空间，具有一个维数$\leq n-2$的闭的无处稠密PL子空间$\Sigma$，其补集为\term{pure}{纯}$n$维定向PL流形。空间$X$带有\term{stratification}{分层}$S$，即由闭PL子空间$X_k$（$0\leq k\leq n$）组成的\term{filtration}{滤过}，满足$X_i\supset X_{i-1}$，$X_i-X_{i-1}$为\term{pure}{纯}$i$维流形，且$(X,S)$沿每个$X_i-X_{i-1}$等奇异。假设$X_{n-1}=X_{n-2}$。对每一个\term{perversity}{反常度}$p$，即满足$p_2=0$且$p_{k+1}=p_k$或$p_k+1$的整数列$p_2,\ldots,p_n$，他们将$X$的第$i$个\term{perversity}{反常度}为$p$的\term{intersection-homology}{交叉同调}群$\mathrm{IH}_i^p(X)$定义为：$X$中与$X_{n-k}$相交部分的维数至多为$i-(k-p_k)$的$i$维奇异PL闭链，模去适当同调关系所得的群。这个群与$S$无关。当$X$正规（即链环连通）且$p=0$（分别地，$p$取最大值）时，他们证明所得的是\term{cohomology}{上同调}群$H^{n-i}(X)$（分别地，同调群$H_i(X)$）。[6]放宽了对$X$的条件，并证明了$\mathrm{IH}_i^p(X)$的拓扑不变性。

我们最关心的情形是：$S$的所有\term{stratum}{层片}都可以取为偶数维，而$p$取中间\term{perversity}{反常度}。按上面的记号，这一\term{perversity}{反常度}由$p_{2k}=k-1$给出（$p_{2k+1}$取$k-1$或$k$均可）。将所得群与$\Q$作张量积后，记为$\mathrm{IH}_i(X)$。它们满足Poincaré\term{duality}{对偶}（[4]、[5]）。将系数从$\Q$\term{extension}{扩张}到$\R$后，可以将它们解释为$L^2$\term{cohomology}{上同调}群，见J. Cheeger在本次会议上的报告。

设$X(\C)$是一个连通正规复代数簇所对应的通常拓扑空间，\SourcePage{15}其复维数为$n$。可以取代数\term{stratification}{分层}$S$，各\term{stratum}{层片}的实维数于是自动为偶数。因此，中间\term{perversity}{反常度}的\term{intersection-homology}{交叉同调}群有定义。若只假设$X$不可约，记其正规化为$X^{\sim}$，我们定义
\[
\mathrm{IH}_i(X(\C)):=\mathrm{IH}_i(X^{\sim}(\C)).
\]

将这些群与$\C$作张量积后，可以用\term{holonomic-module}{完整模}来描述。为给出这一描述，使用代数\term{holonomic-module}{完整模}较为方便；它们是带可积联络并满足适当有限性条件的\term{quasi-coherent}{拟凝聚}模。设$X$嵌入\term{pure}{纯}复$d$维的\term{smooth}{光滑}代数簇$Z$，并假设$X$在$Z$中闭。记$X^\circ$为$X$的\term{smooth}{光滑}部分，$j$为$X^\circ$到$Z$的包含态射。则在唯一同构的意义下，$Z$上存在唯一一个\term{support}{支撑}在$X$上的不可约\term{holonomic-module}{完整模}，记作$j_{!*}\mathcal O_{X^\circ}$；它在$X^\circ$上与\term{holonomic-module}{完整模}$\mathcal H^{d-n}_{X^\circ}(\mathcal O)$相同。若$X$完备，则
\[
\mathrm{IH}_i(X(\C))\otimes\C
=\mathbb H^{d+n-i}\bigl(Z,\Omega^*(j_{!*}\mathcal O_{X^\circ})\bigr).
\]
右边既可以理解为$Z$在Zariski拓扑下、以$\Omega^*(j_{!*}\mathcal O_{X^\circ})$为系数的\term{hypercohomology}{超上同调}，也可以理解为$Z(\C)$在通常拓扑下、以系数为$j_{!*}\mathcal O_{X^\circ}$的全\term{pure}{纯}de Rham\term{complex}{复形}为系数的\term{hypercohomology}{超上同调}。若$X$不必完备，这样得到的是利用局部有限链定义的\term{intersection-homology}{交叉同调}群的变体。$X(\C)$本来的\term{intersection-homology}{交叉同调}群则由$Z(\C)$上以$\Omega^*(j_{!*}\mathcal O_{X^\circ})$为系数的\term{compact-support}{紧支撑}\term{hypercohomology}{超上同调}给出。

不与$\C$作张量积时，$\mathrm{IH}_i(X(\C))$同样可以描述为\term{hypercohomology}{超上同调}群。可以证明，若$X$为$n$维不可约代数簇，则在\term{derived-category}{导出范畴}$\D(X(\C),\Q)$中存在唯一对象$\underline{\mathrm{IC}}$，其\term{cohomology}{上同调}\term{sheaf}{层}$\mathcal H^i(\underline{\mathrm{IC}})$\term{constructible}{可构造}，且当$i\notin[-2n,-n]$时为零，并满足下列三个条件。
\begin{enumerate}[label=(\alph*)]
\item 存在$X$的稠密Zariski开集$U$，使得$\underline{\mathrm{IC}}$在$U(\C)$上同构于仅在$-2n$次放置\term{constant-sheaf}{常值层}$\Q$的\term{complex}{复形}。
\item 当$i>-2n$时，\term{cohomology}{上同调}\term{sheaf}{层}$\mathcal H^i(\underline{\mathrm{IC}})$的\term{support}{支撑}的复余维\SourcePage{16}满足$\geq i-2n+1$。\footnote{译注：此处依原文照录。结合后文条件(b')的\term{equivalence}{等价}写法，此处应为$i+2n+1$。}
\item 对$X$中每个余维$c>0$的不可约子簇$Z$，都存在稠密开集$V\subset Z$，使\term{support}{支撑}\term{cohomology}{上同调}\term{sheaf}{层}$\mathcal H^i_{V(\C)}(\underline{\mathrm{IC}})$在$i\leq-2n+c$时为零。
\end{enumerate}
附带(a)中的同构后，对象$\underline{\mathrm{IC}}$在唯一同构的意义下唯一，并且
\[
\mathrm{IH}_i(X(\C))=\mathbb H_c^{-i}(X(\C),\underline{\mathrm{IC}}).
\]

重新编排$\mathrm{IH}_i(X(\C))$的次数对我们更方便。在本引言余下部分，令
\[
\mathrm{IH}'{}^i(X(\C)):=\mathrm{IH}_{n-i}(X(\C)),
\qquad \underline{\mathrm{IC}}'=\underline{\mathrm{IC}}[-n],
\]
于是
\[
\mathrm{IH}'{}^i(X(\C))=\mathbb H_c^i(X(\C),\underline{\mathrm{IC}}').
\]
$\mathrm{IH}_i(X(\C))$只可能在$i\in[0,2n]$时非零，相应地，$\mathrm{IH}'{}^i(X(\C))$只可能在$i\in[-n,n]$时非零。条件(b)、(c)可改写为：对$X$的每个维数为$d$的真不可约子簇$Z$，都存在稠密Zariski开集$V\subset Z$，使得
\begin{enumerate}[label=\textup{(\alph*')} ,start=2]
\item 当$i\geq-d$时，$\mathcal H^i(\underline{\mathrm{IC}}')$在$V(\C)$上的限制为零；
\item 当$i\leq-d$时，$\mathcal H^i_{V(\C)}(\underline{\mathrm{IC}}')=0$。
\end{enumerate}

设$X$为复代数簇。我们把$X$上的一个\term{perverse-sheaf}{反常层}定义为\term{derived-category}{导出范畴}$\D(X(\C),\Q)$的对象$K$，要求其\term{cohomology}{上同调}\term{sheaf}{层}$\mathcal H^iK$\term{constructible}{可构造}，且在$|i|$充分大时为零（即$K\in\D_c^b(X(\C),\Q)$），并满足以下条件：对$X$中每个维数为$d$的不可约子簇$Z$，都存在稠密Zariski开集$V\subset Z$，使得
\begin{align*}
\text{(b'')}\quad &\mathcal H^i(K)|_{V(\C)}=0 &&(i>-d),\\
\text{(c'')}\quad &\mathcal H^i_{V(\C)}(K)=0 &&(i<-d).
\end{align*}
注意，当$X$不可约时，这里不排除$Z=X$。对于不可约$X$，$\underline{\mathrm{IC}}'$是\term{perverse-sheaf}{反常层}。

条件(b'')\term{equivalence}{等价}于
\[
\dim\Supp\mathcal H^iK\leq-i
\]
（因而当$i>0$时，$\mathcal H^iK=0$）。M. Artin在证明\term{affine}{仿射}簇的\term{cohomology}{上同调}维数定理时，\SourcePage{17}已经引入过这类条件（SGA 4 XIV，第2至4节）。

在代数\term{coherent-sheaf}{凝聚层}理论中，(c'')这类条件很常见，称为深度条件。

\term{perverse-sheaf}{反常层}既不是通常意义的\term{sheaf}{层}，也并不“反常”，因此需要解释这一名称。我们中的一些人并不喜欢“反常”一词。它来自“\term{perversity}{反常度}”：在$\mathrm{IH}^p(X)$中，\term{perversity}{反常度}$p$表示允许闭链偏离与\term{stratification}{分层}横截相交的程度。那么为何称为“\term{sheaf}{层}”？若采用复系数而非有理系数，即把$\D_c^b(X(\C),\Q)$换成$\D_c^b(X(\C),\C)$，且$X$\term{smooth}{光滑}、\term{pure}{纯}$n$维，则函子
\[
M\longmapsto\Omega^*(M)[n]
\]
（解析de Rham\term{complex}{复形}）给出一个\term{equivalence}{等价}，将$X$上具有\term{regular-singularity}{正则奇点}（包括无穷远处）的代数\term{holonomic-module}{完整模}范畴，与$X$上的\term{perverse-sheaf}{反常层}范畴相联系。对于一般的$X$，将它作为闭子簇嵌入\term{pure}{纯}$d$维\term{smooth}{光滑}簇$Z$后，应考察$Z$上\term{support}{支撑}在$X$中的\term{holonomic-module}{完整模}，以及函子
\[
M\longmapsto\Omega^*(M)[d]|_{X(\C)}.
\]
\term{smooth}{光滑}$X$上的\term{holonomic-module}{完整模}本身是\term{sheaf}{层}（$\mathcal D_X$-模\term{sheaf}{层}），它们的性质启发了\term{perverse-sheaf}{反常层}的若干性质。例如有如下结论；下面重新取$\Q$为系数域。
\begin{itemize}
\item $X$上的\term{perverse-sheaf}{反常层}构成\term{abelian-category}{阿贝尔范畴}。
\item 它们是局部性质的对象：$X$的各个Zariski开集$U$上的\term{perverse-sheaf}{反常层}范畴构成一个\term{stack}{栈}。需要注意，\term{derived-category}{导出范畴}$\D_c^b(U(\C),\Q)$并非如此。
\item 设$j:U\hookrightarrow X$为开嵌入，$i:F\hookrightarrow X$为其闭补集的嵌入，则存在与通常\term{sheaf}{层}的形式体系相近的一套形式体系，具有函子$j_!,j^*,j_*,i^*,i_*,i^!$。对于这些函子在\term{perverse-sheaf}{反常层}上的对应物，我们在左上角标以$p$。有一条经验法则，说明关于\term{sheaf}{层}范畴的哪些结果可以推广到\term{perverse-sheaf}{反常层}：将\term{sheaf}{层}范畴换成其反范畴，并交换$j_!$与$j_*$、$i^*$与$i^!$后得到的\term{duality}{对偶}陈述，也必须成立。例如：
\end{itemize}
\begin{enumerate}[label=\alph*)]
\item 对$X$上的\term{perverse-sheaf}{反常层}$F$，有\term{exact-sequence}{正合列}
\[
0\longrightarrow{}^pi_*{}^pi^!F\longrightarrow F
\longrightarrow{}^pj_*{}^pj^*F.
\]
\SourcePage{18}
这对通常\term{sheaf}{层}也成立，其\term{duality}{对偶}为\term{exact-sequence}{正合列}
\[
j_!j^*F\longrightarrow F\longrightarrow i_*i^*F\longrightarrow0.
\]
但对\term{perverse-sheaf}{反常层}而言，不能像通常\term{sheaf}{层}那样在这个\term{duality}{对偶}序列的左端再添一个零。
\item 对$U$上的\term{perverse-sheaf}{反常层}$F$，存在自然态射
\[
\alpha:{}^pj_!F\longrightarrow{}^pj_*F.
\]
对通常\term{sheaf}{层}来说，它是单态射。但\term{duality}{对偶}陈述“它是满态射”一般不成立；对于\term{perverse-sheaf}{反常层}，$\alpha$一般也不是单态射。函子
\[
j_{!*}:F\longmapsto\im({}^pj_!F\longrightarrow{}^pj_*F)
\]
称为\term{intermediate-extension}{中间延拓}，对我们极为有用。
\end{enumerate}
若$X$不可约、维数为$n$，$U$为\term{smooth}{光滑}稠密Zariski开集，而$\Q[n]$是$U(\C)$上由$-n$次的\term{constant-sheaf}{常值层}$\Q$构成的\term{perverse-sheaf}{反常层}，则
\[
\underline{\mathrm{IC}}'(X)=j_{!*}(\Q[n]).
\]

\term{perverse-sheaf}{反常层}在某些方面比通常\term{sheaf}{层}具有更好的性质。例如：
\begin{itemize}
\item 在\term{perverse-sheaf}{反常层}的\term{abelian-category}{阿贝尔范畴}中，每个对象均为有限长度（参见\term{holonomic-module}{完整模}的对应陈述）。
\item 若$j:U\hookrightarrow X$为\term{affine}{仿射}开集的包含，则函子${}^pj_*$正合。
\item $\D_c^b(X(\C),\Q)$上的Verdier\term{duality}{对偶}将\term{perverse-sheaf}{反常层}变为\term{perverse-sheaf}{反常层}。这解释了\term{intersection-homology}{交叉同调}中的Poincaré\term{duality}{对偶}；它也与\term{smooth}{光滑}$X$上的\term{holonomic-module}{完整模}具有良好\term{duality}{对偶}理论这一事实相呼应。
\end{itemize}
通常的\term{constructible}{可构造}\term{sheaf}{层}与\term{perverse-sheaf}{反常层}，各自都构成$\D_c^b(X(\C),\Q)$的一个\term{abelian-subcategory}{阿贝尔子范畴}$C$。在这两种情形下，都有从$\D_c^b(X(\C),\Q)$到$C$的函子$H^i$（“\term{cohomology}{上同调}\term{sheaf}{层}”），服从同样的形式体系：长\term{exact-sequence}{正合列}、把$K$逐步分解为各个$H^iK[-i]$等。通常\term{sheaf}{层}的范畴$C$适合于\term{inverse-image}{逆像}函子，\term{perverse-sheaf}{反常层}的范畴则适合于Verdier\term{duality}{对偶}，以及某些\term{direct-image}{正像}的研究（\term{affine}{仿射}态射、\term{vanishing-cycles}{消失轮换}、相对\term{hard-lefschetz}{强Lefschetz定理}等）。

\SourcePage{19}
对于代数闭域$k$上的\term{finite-type}{有限型}\term{scheme}{概形}$X$，\term{perverse-sheaf}{反常层}的形式体系可以直接移植到\term{l-adic}{$\ell$-进}\term{cohomology}{上同调}中，其中$\ell$与$k$的特征指数互素。只需把$\D_c^b(X(\C),\Q)$换成$\Q_\ell$-线性\term{derived-category}{导出范畴}$\D_c^b(X,\Q_\ell)$。特别地，若$X$不可约，就可以定义\term{intersection-homology}{交叉同调}群$\mathrm{IH}_i(X,\Q_\ell)$。当$k=\C$时，这些群与$\mathrm{IH}_i(X(\C))\otimes\Q_\ell$典范同构。

对任意基域，障碍在于文献中尚无可用的$\D_c^b(X,\Q_\ell)$定义。O. Gabber已有一个定义，但尚未写出。所幸在我们尤其关心的有限域情形，一个朴素的定义就已足够。其所以足够，是因为对每一个有限域，都有$H^*(\Gal(\bar k/k),\Z/\ell)$有限。

主要结果的一种表述是：若$X_0$为有限域$k$上的\term{finite-type}{有限型}\term{scheme}{概形}，则$X_0$上的每一个单\term{mixed}{混合}\term{perverse-sheaf}{反常层}$F_0$都是\term{pure}{纯}的。“\term{mixed}{混合}”这一性质在所有通常运算下保持，故在实际应用中几乎总能满足；“单”是指\term{perverse-sheaf}{反常层}\term{abelian-category}{阿贝尔范畴}中的\term{simple-object}{单对象}。对于不可约的$X_0$，\term{perverse-sheaf}{反常层}$\underline{\mathrm{IC}}'$既单又\term{mixed}{混合}。记$q$为$k$的元素个数，$\bar k$为$k$的代数闭包，$X=X_0\otimes_k\bar k$，$F$为$F_0$在$X$上的\term{inverse-image}{逆像}。“\term{pure}{纯}”是一个局部性质。当$X$完备时，它保证存在一个整数$w$（即$F_0$的\term{weight}{权}），使\term{frobenius}{Frobenius} $F^*$在$H^i(X,F)$上的特征值均为代数数，且其所有复共轭的绝对值都是$q^{(w+i)/2}$。

在一维情形，这个结果\term{equivalence}{等价}于[1]的1.8.4。后来，D. Kazhdan与G. Lusztig在与Schubert簇有关的若干特殊情形中得到了它。第一个一般性证明由O. Gabber给出，与本书的证明不同。该定理提供了大量\term{pure}{纯}\term{complex}{复形}，从而使[1]第6节对\term{pure}{纯}\term{complex}{复形}的研究获得了实质内容。

下面概述本文各章。

\par\medskip\noindent\textbf{第1章。} \term{triangulated-category}{三角范畴}方面，我们主要以J. L. Verdier的[10]为参考。1.1节给出补充说明、动机和例子。
\SourcePage{20}
1.2节研究带有满子范畴$C$的\term{triangulated-category}{三角范畴}$\mathcal D$，要求满足公理1.2.0及1.2.4(ii)，即所谓\term{admissible}{容许}\term{abelian-subcategory}{阿贝尔子范畴}。设$A$是具有足够多\term{injective}{内射}对象的\term{abelian-category}{阿贝尔范畴}，$\mathcal D=\D^+(A)$，并将其视为$A$中\term{injective}{内射}对象组成的下有界\term{complex}{复形}范畴，其中态射取到\term{homotopy}{同伦}为止。则公理1.2.0保证：若$X,Y\in C$，且\term{complex}{复形}态射$f,g:X\rightrightarrows Y$\term{homotopy}{同伦}，则这一\term{homotopy}{同伦}在二阶\term{homotopy}{同伦}的意义下唯一。也就是说，$f$与$g$之间的两个\term{homotopy}{同伦}$H_1,H_2$由一个二阶\term{homotopy}{同伦}$H'$联系起来；而$H'$又在三阶\term{homotopy}{同伦}的意义下唯一，如此类推。我们证明$C$是\term{abelian-category}{阿贝尔范畴}。

$\mathcal D$上的每一个\term{t-structure}{$t$-结构}（1.3.1）都确定了$\mathcal D$的一个\term{admissible}{容许}\term{abelian-subcategory}{阿贝尔子范畴}$C$（1.2.5、1.3.6），称为该\term{t-structure}{$t$-结构}的\term{heart}{心}；反方向也有一个部分成立的结论（1.3.13）。设$A$为\term{abelian-category}{阿贝尔范畴}，$C$为$A$的严格满\term{abelian-subcategory}{阿贝尔子范畴}，对取核、余核及\term{extension}{扩张}封闭；记$\D_C(A)$为\term{derived-category}{导出范畴}中由所有满足$H^iK\in C$的$K$组成的满子范畴。则满子范畴对
\[
\bigl(\D_C^{\leq0}(A),\D_C^{\geq0}(A)\bigr)
\]
给出$\D_C(A)$上的一个\term{t-structure}{$t$-结构}，其\term{heart}{心}为$C$。1.3节展示了一条经验法则：在这一背\term{site}{景}下成立的结论，对所有\term{t-structure}{$t$-结构}也成立。例如，有\term{truncation}{截断}函子$\tau$和\term{cohomological-functor}{上同调函子}$H^0:\mathcal D\to C$等。不过有一个例外：在3.1节中，为构造自然函子$\D^b(C)\to\mathcal D$，我们需要用到$\mathcal D$上的额外数据。

我们还研究$t$-范畴（即带\term{t-structure}{$t$-结构}的\term{triangulated-category}{三角范畴}）之间的\term{exact-functor}{正合函子}$T$如何在各自的\term{heart}{心}之间诱导函子${}^pT$，并研究它的正合性和\term{adjunction}{伴随}性质。

1.4节给出一种构造\term{t-structure}{$t$-结构}的方法。第2章将用它来构造以\term{perverse-sheaf}{反常层}范畴为\term{heart}{心}的\term{t-structure}{$t$-结构}。粗略地说，给定拓扑空间$X$、开集嵌入$j:U\hookrightarrow X$及其闭补集的嵌入$i:F\hookrightarrow X$，我们要从$\D(U)$与$\D(F)$上的\term{t-structure}{$t$-结构}构造$\D(X)$上的\term{t-structure}{$t$-结构}。即使$\D(U)$与$\D(F)$只带有各自自然\term{t-structure}{$t$-结构}（1.3.2）的平移，所得的\term{t-structure}{$t$-结构}也已不再显然。

\term{l-adic}{$\ell$-进}\term{cohomology}{上同调}的需要迫使我们\SourcePage{21}采用更抽象的框架，把$\D(X)$、$\D(U)$和$\D(F)$之间的关系公理化。我们就在这一框架中建立\term{perverse-sheaf}{反常层}范畴之间的函子$j_!,j^*,j_*,i^*,i_*,i^!$的形式体系，以及\term{intermediate-extension}{中间延拓}$j_{!*}$的形式体系。

\par\medskip\noindent\textbf{第2章。} 在2.1节中，我们考虑\term{stratification}{分层}空间$(X,S)$，迭代应用1.4节的构造，为每一个\term{perversity}{反常度}函数$p:S\to\Z$赋予$\D(X)$上的\term{t-structure}{$t$-结构}。具体做法是把各个$\D(S)$（这里$S$为一个\term{stratum}{层片}）的自然\term{t-structure}{$t$-结构}平移$p(S)$后加以\term{recollement}{粘合}。2.2节转向\term{scheme}{概形}；通过对越来越精细的代数\term{stratification}{分层}取极限，从2.1节的\term{t-structure}{$t$-结构}得到相应的\term{t-structure}{$t$-结构}。我们也讨论\term{l-adic}{$\ell$-进}情形。

\par\medskip\noindent\textbf{第3章。} 本章给出若干技术性较强的补充内容。初次阅读时可以跳过。

3.1节使用\term{filtered-derived-category}{滤过导出范畴}。在某些情形下，对\term{heart}{心}为$C$的$t$-范畴$\mathcal D$，它们使我们能够定义函子$\D^b(C)\to\mathcal D$（3.1.10）。这个函子\term{fully-faithful}{全忠实}，当且仅当对每个$i>0$，函子
\[
C^{\circ}\times C\longrightarrow(\mathrm{Ab}),\qquad
(X,Y)\longmapsto\Hom_{\mathcal D}(X,Y[i])
\]
可\term{effaceable}{消去}（3.1.16）。

3.2节给出将\term{derived-category}{导出范畴}中局部给定的对象\term{recollement}{粘合}起来的判据。在$i<0$时$\mathcal H^i\mathrm{RHom}$\term{sheaf}{层}为零的假设下，\term{recollement}{粘合}可以像通常\term{sheaf}{层}的\term{recollement}{粘合}那样顺利进行。这适用于\term{perverse-sheaf}{反常层}。\term{l-adic}{$\ell$-进}情形已在2.2.19中用另一种方法处理。

3.3节说明，当系数取$\Z$（或$\Z_\ell$）而不是$\Q$（或$\Q_\ell$）时，应如何调整\term{intersection-cohomology}{交叉上同调}的形式体系。

在第4、5章中，除4.4节外，我们只考虑域$k$上的\term{finite-type}{有限型}\term{scheme}{概形}，并且只取中间\term{perversity}{反常度}。第4章给出几何结果：即使假设$k$代数闭，一般性也几乎没有损失。第5章中，$k$为有限域，研究的现象具有算术性质；\term{extension}{扩张}到$k$的代数闭包后，它们又产生几何上的推论。

\SourcePage{22}
\par\medskip\noindent\textbf{第4章。} 4.1节在\term{perverse-sheaf}{反常层}这一自然框架中，陈述M. Artin关于\term{affine}{仿射}态射的高阶\term{direct-image}{正像}的定理（SGA 4 XIV 3.1）。4.4节对“\term{nearby-cycles}{邻近轮换}”函子证明类似定理。

4.2节以\term{perverse-sheaf}{反常层}及其函子性的语言，给出4.1节及其他已知结果的若干推论。

4.3节证明\term{perverse-sheaf}{反常层}范畴既是\term{artinian}{阿廷}的，也是\term{noetherian}{诺特}的。

4.5节为\term{pure}{纯}性定理5.3.4的证明作准备。它可以看作对“中间维数的\term{cohomology}{上同调}群最大、也最有意义”这一主题的一种展开。

\par\medskip\noindent\textbf{第5章。} 设$X_0$为有限域$k$上的分离\term{finite-type}{有限型}\term{scheme}{概形}，$\bar k$为$k$的代数闭包，$X$为$X_0$经系数域\term{extension}{扩张}所得的$\bar k$-\term{scheme}{概形}。5.1节回顾[1]中有关\term{weight}{权}的形式体系的若干定义和结果，还研究$X_0$上的\term{perverse-sheaf}{反常层}与$X$上带\term{frobenius}{Frobenius}作用的\term{perverse-sheaf}{反常层}之间的关系。

由[1]的3.3.1、6.2.3可知，若$K_0\in\D_c^b(X_0,\overline{\Q}_\ell)$的\term{weight}{权}$\leq w$，则$H_c^i(X,K)$的\term{weight}{权}$\leq w+i$。对于\term{perverse-sheaf}{反常层}$K_0$，定理5.2.1给出这一结论的逆命题；更准确地说，是其\term{duality}{对偶}结论的逆命题，即把$H_c^i$换成$H^i$，把$\leq$换成$\geq$。我们需要检验的是$H^0(U,K)$，其中$U$为$X$上的\term{affine}{仿射}\term{etale}{平展}\term{scheme}{概形}。应当把这个$H^0$看作中间维数的\term{cohomology}{上同调}群。对\term{affine}{仿射}$U$，由4.1.1可知，函子$K\mapsto H^0(U,K)$（$K$为\term{perverse-sheaf}{反常层}）右正合。它的作用可以与通常\term{sheaf}{层}的“在一点处取纤维”函子相比（参见4.1.6）。

5.3节由5.2.1推出：\term{weight}{权}$\leq w$的\term{mixed}{混合}\term{perverse-sheaf}{反常层}的任意\term{subquotient}{子商}仍然\term{weight}{权}$\leq w$（5.3.1）；\term{intermediate-extension}{中间延拓}保持\term{weight}{权}（5.3.2、5.3.3）；单\term{mixed}{混合}\term{perverse-sheaf}{反常层}是\term{pure}{纯}的（5.3.4）。若$K_0,L_0$为$X_0$上的两个\term{perverse-sheaf}{反常层}，则在\term{perverse-sheaf}{反常层}\term{abelian-category}{阿贝尔范畴}中$K_0$被$L_0$\term{extension}{扩张}的\term{extension}{扩张}类群，等于在$\D_c^b(X_0)$中计算的$\Hom(K_0,L_0[1])$。利用\term{weight}{权}不为零的\term{frobenius}{Frobenius}特征值不可能\SourcePage{23}等于$1$，可以证明某些$\Ext^1$群为零。由此保证每个\term{mixed}{混合}\term{perverse-sheaf}{反常层}上存在一个递增的\term{weight-filtration}{权滤过}$W$，且$\gr_i^W$\term{pure}{纯}、\term{weight}{权}为$i$（5.3.5）；这一\term{filtration}{滤过}还具有函子性。

设$K_0\in\D_c^b(X_0,\overline{\Q}_\ell)$，记$\mathcal H^iK_0$为其通常意义的\term{cohomology}{上同调}\term{sheaf}{层}（$\overline{\Q}_\ell$-\term{sheaf}{层}），${}^p\mathcal H^iK_0$为其反常意义的\term{cohomology}{上同调}\term{sheaf}{层}（\term{perverse-sheaf}{反常层}）。假设$K_0$\term{mixed}{混合}。按定义，$K_0$的\term{weight}{权}$\leq w$，当且仅当每个$\mathcal H^iK_0$的\term{weight}{权}$\leq w+i$。在5.4.1中，我们证明这也\term{equivalence}{等价}于：每个${}^p\mathcal H^i(K_0)$的\term{weight}{权}$\leq w+i$。由\term{duality}{对偶}性，“\term{weight}{权}$\geq w$”也可以在${}^p\mathcal H^i(K_0)$上检验，而不能在$\mathcal H^iK_0$上检验。

5.4.5与5.3.8给出两个\term{semisimple}{半单}性结果。两者都从$X_0$上的\term{pure}{纯}\term{complex}{复形}$K_0$出发，而结论是关于其在$X$上的\term{inverse-image}{逆像}$K$的：假设是算术的，结论是几何的。5.4.5证明$K$同构于各个${}^p\mathcal H^i(K)[-i]$的直和。5.3.8证明每个${}^p\mathcal H^iK$都是单\term{perverse-sheaf}{反常层}的直和。再结合\term{pure}{纯}性在\term{proper}{固有}态射的\term{direct-image}{正像}下保持这一事实，可以推出局部及整体\term{invariant-cycles}{不变轮换}定理的推广（5.4.7、5.4.8，以及5.4.9的加强）。5.4.10证明\term{hard-lefschetz}{强Lefschetz定理}的相对版本。

\par\medskip\noindent\textbf{第6章。} 第5章的结果对有限域的代数闭包$\F$上的\term{finite-type}{有限型}\term{scheme}{概形}的\term{cohomology}{上同调}具有几何上的推论。6.1节说明如何将这些结果从$\F$转移到$\C$。如果所讨论的是代数闭域理论的一阶陈述，只需指出$\C$同构于一族有限域代数闭包的超积即可。但是，我们关心的陈述至少乍看之下并不\term{equivalence}{等价}于一阶陈述，因而必须利用\term{constructible}{可构造}性等结果；这些结果可以理解为表达了\term{l-adic}{$\ell$-进}\term{cohomology}{上同调}的某种可计算性。

6.2节应用这些原理。主要结果是6.2.5、6.2.8、6.2.9及6.2.10。

还应指出，以下内容本应收入本书，但我们未能完成。
\SourcePage{24}
\begin{itemize}
\item \term{perverse-sheaf}{反常层}与\term{holonomic-module}{完整模}的关系。如本引言所述，这一关系起到了重要的启发作用。主要陈述是4.1.9，本书未给出证明。
\item \term{vanishing-cycles}{消失轮换}函子（4.4节的一个部分结果除外）。假定已经有可用的\term{l-adic}{$\ell$-进}\term{derived-category}{导出范畴}定义，并已定义那些“显然”的态射——也就是说，仍有不少工作要做——则下列陈述的证明是已知的。
\end{itemize}
\begin{enumerate}[label=\alph*)]
\item 设$X$为离散赋值环的谱$(S,\eta,s)$上的\term{finite-type}{有限型}\term{scheme}{概形}，
\[
\mathrm R\Psi:\D_c^b(X,\Q_\ell)
\longrightarrow\D_c^b(X\times_S\bar s,\Q_\ell)
\]
为“\term{nearby-cycles}{邻近轮换}”函子。它是\term{t-exact}{$t$-正合}的，特别地，它在\term{perverse-sheaf}{反常层}范畴之间诱导\term{exact-functor}{正合函子}，并与Verdier\term{duality}{对偶}交换。
\item 若$S$为有限域上一条曲线$S_1$在某闭点处的Hensel化，$X/S$由$S_1$上的\term{finite-type}{有限型}\term{scheme}{概形}$X_1$经\term{base-change}{基变换}得到，而$F_1$为$X_1$上的\term{mixed}{混合}\term{perverse-sheaf}{反常层}，则$F_1$与$\mathrm R\Psi_\eta(F_1|_{X_\eta})$的\term{weight-filtration}{权滤过}和\term{monodromy}{单值作用}之间存在一个关系，推广了[1]的1.8.5。
\end{enumerate}
6.1节也未说明如何将(b)的几何推论从$\F$转移到$\C$。

最后，我们没有讨论Fourier变换。

\section*{0.0. 记号与术语}
\addcontentsline{toc}{section}{0.0. 记号与术语}
读者可以在文末找到术语索引和符号索引，列有本书使用的主要新术语、新记号或非标准的术语与记号。

注意，从1.4节起，我们通常将\term{sheaf}{层}范畴的\term{derived-category}{导出范畴}之间的函子简记为$f_*,f^*,f_!,f^!$；它们通常记作$\mathrm Rf_*,\mathrm Rf^*$（或$\mathrm Lf^*$）、$\mathrm Rf_!$及$\mathrm Rf^!$。通常\term{sheaf}{层}范畴之间的同名函子则在左上角标以$0$，因为它们对应于\term{perversity}{反常度}$0$。
